% ECONOMICS_PROBLEM: {"id": "C62-48", "title": "Sequential and recursive equilibrium with a distributional state", "jel_code": "C62", "source_id": "OP-0583"}
% FIRST_POSTED: 2026-10-09
% ATTEMPT_STATUS: unresolved
% ENTRY_KIND: research_attempt
% !TeX program = pdflatex
% Self-contained: no external input or bibliography files required.
\documentclass[11pt,a4paper]{article}
\usepackage[T1]{fontenc}
\usepackage[utf8]{inputenc}
\usepackage{lmodern}
\usepackage[margin=24mm]{geometry}
\usepackage{amsmath,amssymb,amsthm,mathtools}
\usepackage{booktabs,longtable,array,enumitem,microtype,xurl,hyperref}
\hypersetup{colorlinks=true,linkcolor=blue,citecolor=blue,urlcolor=blue,pdftitle={Checked research attempts: nine economics and stochastic-analysis problems}}
\setlength{\parindent}{0pt}
\setlength{\parskip}{3pt}
\setlength{\emergencystretch}{3em}
\widowpenalty=10000\clubpenalty=10000\displaywidowpenalty=10000
\setlist[enumerate]{leftmargin=*,itemsep=2pt,topsep=3pt}
\newcommand{\R}{\mathbb R}
\newcommand{\N}{\mathbb N}
\newcommand{\E}{\mathbb E}
\newcommand{\Pp}{\mathbb P}
\newcommand{\Q}{\mathbb Q}
\newcommand{\F}{\mathcal F}
\newcommand{\ind}{\mathbf 1}
\newcommand{\cE}{\mathcal E}
\newcommand{\Law}{\operatorname{Law}}
\newcommand{\sgn}{\operatorname{sgn}}
\newcommand{\Var}{\operatorname{Var}}
\newcommand{\dist}{\operatorname{dist}}
\newcommand{\KL}{\operatorname{KL}}
\newcommand{\TV}{\operatorname{TV}}
\newcommand{\Lip}{\operatorname{Lip}}
\DeclareMathOperator*{\esssup}{ess\,sup}
\DeclareMathOperator*{\argmax}{arg\,max}
\newtheorem{theorem}{Theorem}
\newtheorem{lemma}[theorem]{Lemma}
\newtheorem{proposition}[theorem]{Proposition}
\newtheorem{corollary}[theorem]{Corollary}
\theoremstyle{definition}
\newtheorem{example}[theorem]{Example}
\newcommand{\report}[3]{\clearpage\setcounter{theorem}{0}\renewcommand{\thetheorem}{#1.\arabic{theorem}}\renewcommand{\theHtheorem}{#1.\arabic{theorem}}\section*{#1: #2}\addcontentsline{toc}{section}{#1: #2}\textbf{Outcome:} #3.\par\textbf{Tools:} web browsing available; Python computation available. Literature checked on 9 October 2026.}
\newcommand{\stage}[2]{\subsection*{#1) #2}}
\newcommand{\unresolved}{\textbf{LABEL: UNRESOLVED.}}
\newcommand{\disproof}{\textbf{LABEL: FULL SOLUTION. SUBLABEL: COUNTEREXAMPLE/DISPROOF.}}

\begin{document}
\section*{C62-48: Sequential versus reduced-state recursive equilibrium}
\textbf{Outcome:} UNRESOLVED as a general programme; a conditional equivalence theorem and a selection obstruction proved.\par\textbf{Tools:} web browsing available; Python computation available. Literature checked on 9 October 2026.
\subsection*{1) FORMAL RESTATEMENT}
The uploaded question asks which sequential equilibria of the economy called ``MA-01'' admit representation on the aggregate state $s=(z,\mu)$ and household state $(a,e)$, with prices, policies, continuation values, and the distribution transition all measurable functions of these states. The state $\mu$ is the full cross-sectional distribution, not a finite list of moments.

\textbf{Missing specification.} The nine supplied standalone files do not define MA-01. The minimal correction is to supply utility, production, shock kernels, constraints, initial distributions and transversality, together with the continuation domain. Below we state a bounded Markov specification, not a reconstruction of MA-01.

\textbf{Quantified representation property.} For each relevant aggregate history $h$, let $s(h)=(z(h),\mu(h))$. Let $P_h$ denote the equilibrium price vector and $g_h(a,e)$ the savings policy, defined for every individual state under consideration. A representation requires measurable $P,g,V,\Phi$ on an invariant domain $D$ such that, for every relevant history and every relevant individual state,
\[
 P_h=P(s(h)),\quad g_h(a,e)=g(a,e,s(h)),\quad
 V_h(a,e)=V(a,e,s(h)),
\]
and $s(hz')=\Phi(s(h),z')$ for every possible successor shock. Equalities for policies may be imposed modulo the explicitly chosen individual-state null sets; in our theorem they hold pointwise. Invariance means that every positive-probability successor remains in $D$. A representation on one realized shock path is insufficient.

\subsection*{2) QUICK LITERATURE/CONTEXT CHECK}
Cao's primary working paper \cite{Cao} constructs a generalized recursive equilibrium in a specified Krusell--Smith environment. This supports separating state construction and equilibrium existence from the smaller-state representation question. The uploaded historical citations do not themselves identify the missing MA-01 primitives. We therefore do not assert either universal equivalence or its universal failure for that unspecified model.

\subsection*{3) ATTACK PLAN}
\textbf{Proof routes.} Prove Bellman verification on a fully specified invariant state space; derive a fibre-consistency criterion for a given sequential equilibrium; impose uniqueness of individual maximizers as a sufficient way of eliminating policy-selection dependence on history.

\textbf{Disproof routes.} Look for two histories with the same $(z,\mu)$ and different prices; alternatively exploit ties in optimal saving choices while keeping the distribution unchanged; finally test whether unbounded tails invalidate the transversality step. The clean construction below uses the second route.

\textbf{Landscape and tools.} Bellman contraction identifies values; measurable selection defines policies; pushforwards express distribution consistency; fibre factorization tests collisions; conditional iteration verifies deviations; discounted tail bounds justify limits; and telescoping budgets construct the obstruction.

\subsection*{4) WORK}
\paragraph{Explicit bounded specification for the theorem.}
Let the asset interval $A$ be compact, and let the efficiency set $E$ and aggregate-shock set $Z$ be finite. Set $X=A\times E$, and let $D\subset Z\times\mathcal P(X)$ be a Borel invariant aggregate-state domain. Time is $t\in\N=\{0,1,2,\ldots\}$, with discount factor $0<\beta<1$. At state $(x,s)$ the household chooses an action $a'$ in a nonempty feasible set $\Gamma(x,s)$, receiving a measurable reward $\ell(x,s,a')$ with $|\ell|\le M$. Given Borel price functions $P$ and an aggregate transition $\Phi$, a specified Markov probability kernel $Q(dx',ds'\mid x,s,a')$ describes the next individual and aggregate states. The individual is a price taker: a unilateral deviation does not change $P$ or $\Phi$.

Assume the Bellman operator
\[
 (\mathcal BV)(x,s)=\sup_{a'\in\Gamma(x,s)}
 \left\{\ell(x,s,a')+\beta\int V(x',s')Q(dx',ds'\mid x,s,a')\right\}
\]
maps bounded Borel functions to bounded Borel functions and its fixed point has a Borel maximizing selector. These are explicit measurability/attainment assumptions, satisfied in appropriate compact continuous control specifications but not automatic for arbitrary Borel primitives. Define market clearing, firm optimality, and resource feasibility as part of the declared economy. Distribution consistency for deterministic saving policy $g$ is, for each possible $z'$,
\[
 \mu'(da',de')=\int_{A\times E}\delta_{g(a,e,s)}(da')
        P_E(de'\mid e,z,z')\,\mu(da,de),
\]
where $P_E$ is the declared efficiency kernel. Both equilibrium notions in this theorem are distributional: cross-sectional states are conditional laws satisfying this equality. No additional agent-indexed continuum implementation or unproved exact law of large numbers is invoked.

\begin{theorem}[Recursive verification and a conditional converse]
Under the bounded specification, the following statements hold.

(a) Suppose Borel $P,g,V,\Phi$ satisfy Bellman optimality with $g$ maximizing, distribution consistency, and all declared market and firm conditions on $D$. Then for every initial $s_0\in D$ they induce a sequential competitive equilibrium, optimal against all admissible history-dependent household deviations.

(b) Fix one initial aggregate state and all of its relevant finite aggregate histories. Suppose a sequential equilibrium specifies optimal household continuations for every individual state at every such history. Suppose its prices, policies, and successor aggregate states are constant on fibres of the state map $h\mapsto s(h)$, and the primitive household feasible sets, rewards, and shock kernels depend only on the stated Markov states. Then it has a recursive representation on the invariant set of all these reachable aggregate states. Its continuation values necessarily factor through the same states. This converse does not assert existence of such a fibre-consistent sequential equilibrium.
\end{theorem}
\begin{proof}
For any bounded $V,W$, taking the difference of the two suprema gives
$\|\mathcal BV-\mathcal BW\|_\infty\le\beta\|V-W\|_\infty$.
For completeness, starting from zero the successive Bellman iterates have successive differences at most $M\beta^j$. The geometric series makes them uniformly Cauchy, their limit is bounded Borel, and the displayed Lipschitz inequality makes the limit a fixed point. If two fixed points existed, their distance would be at most $\beta$ times itself, so they would coincide. The fixed point satisfies $\|V\|_\infty\le M/(1-\beta)$.

For (a), use the kernels and policy to define the sequential laws starting from any $s_0\in D$. Existence of a process with these laws is the usual iterated-kernel construction on standard Borel spaces; its hypotheses hold because the domains and kernels here are Borel. Bellman optimality implies, under any admissible deviation,
\[
 V(X_t,S_t)\ge \ell(X_t,S_t,a'_{t})
       +\beta\E[V(X_{t+1},S_{t+1})\mid\text{the household history at }t].
\]
This conditional inequality holds even if the chosen action uses additional past information, since it holds for every feasible action at the current state and the transition is Markov conditional on that action. Iterating and taking expectations gives
\[
 V(x_0,s_0)\ge\E\sum_{t=0}^{N-1}\beta^t\ell_t
                   +\beta^N\E V(X_N,S_N).
\]
The last term has absolute value at most $\beta^NM/(1-\beta)$ and tends to zero. The discounted reward series converges absolutely, dominated by $M/(1-\beta)$. Under $g$ every Bellman inequality is equality, so $g$ attains the upper bound and is optimal against every such deviation. By the declared pushforward equality and invariance, all subsequent distributions and aggregate states agree with the recursive construction. The stipulated market and firm conditions hold at every reached state. This proves (a).

For (b), finite $Z$ and one initial state give a countable set of finite aggregate histories and hence a countable reachable state set $D_0$. Include every positive-probability successor, not merely a single realized path. Fibre consistency defines $P(s)$, $g(x,s)$, and $\Phi(s,z')$ independently of the chosen representative history. Enumerate the histories and use the first representative of each state. Each $g(\cdot,s)$ is measurable by the sequential specification, and countability of $D_0$ makes the combined function Borel on $X\times D_0$: preimages are countable unions of its measurable slices times singleton states. Prices and transitions are measurable for the same reason.

The resulting fixed-price, fixed-transition household control problem is the same at any two histories in a fibre: feasible actions, rewards, and conditional transition laws coincide at every subsequent step. The Bellman argument already proved supplies its unique bounded optimal value $V$. Sequential optimality for every individual state therefore identifies each $V_h(x)$ with $V(x,s(h))$. Moreover its policy must attain the Bellman maximum. Otherwise at some current state choosing a maximizing action and then its Borel optimal continuation would strictly improve upon the declared sequential continuation, contradicting optimality. Distribution consistency and markets transfer directly from all the sequential histories. This proves the representation on $D_0$.
\end{proof}

\paragraph{What the fibre condition says, and does not say.}
For a given sequential equilibrium, any state representation necessarily gives identical prices and policies whenever both the aggregate and individual states coincide. The theorem makes that necessary condition sufficient on the stated countable reachable domain when the other Markov primitives and continuation requirements hold. On an uncountable prescribed domain of initial states, mere set-theoretic constancy on fibres is not enough: Borel factor maps, or a separate measurable-factorization theorem with verified hypotheses, are also needed.

If prices and aggregate transitions already factor and each fixed-state Bellman maximizer is unique, sequential optimality forces policy fibre consistency. This is a useful sufficient condition, but strict concavity of utility alone is not a proof that equilibrium prices or aggregate transitions factor. Nor does the theorem construct an equilibrium or prove uniqueness of the aggregate equilibrium.

\begin{example}[Same full distribution, different optimal policies]
Consider the following fully specified auxiliary economy, not asserted to equal the missing MA-01. There is a unit continuum of households, one deterministic efficiency and aggregate state, labor supply one per household, $A=[0,1]$, $\beta=1/2$, and $u(c)=c$. Production is
$F(K,L)=2K+2L$, with full capital depreciation. Competitive gross capital return and wage are both $2$; firms earn zero profit at these prices and any demanded nonnegative inputs maximize profit. Household budgets are
\[
 c_t=2a_t+2-a_{t+1},\qquad a_{t+1}\in[0,1].
\]
Every such asset path is feasible, with $1\le c_t\le4$. For every $N$,
\[
 \sum_{t=0}^{N-1}2^{-t}c_t
 =2a_0+4(1-2^{-N})-2^{-(N-1)}a_N.
\]
This follows by expanding the budget sum: each interior asset appears with coefficients $2\cdot2^{-t}$ and $-2^{-(t-1)}$, which cancel. Bounded assets make the last term vanish, so every feasible path has utility $2a_0+4$ and is optimal.

Initialize the full distribution at $\mu=(\delta_0+\delta_1)/2$. At even dates use $g_t(a)=a$; at odd dates use $g_t(a)=1-a$. Both maps preserve $\mu$. Thus $K_t=K_{t+1}=1/2$ at all dates. Aggregate consumption is
$2K_t+2-K_{t+1}=5/2$, gross investment is $1/2$, and output is $3$, so goods and factor markets clear. The preceding household argument proves sequential optimality, including off-support individual states if their policies are extended by the same formulas.

The aggregate state $(z,\mu)$ never changes. At an individual state $a=0$, which has mass $1/2$ at every date, saving is $0$ at even dates and $1$ at odd dates. No single time-independent function $g(a,e,z,\mu)$ can equal both. Hence this sequential equilibrium has no policy representation on the proposed reduced state. Its value nevertheless does factor: $V(a)=2a+4$. Adding that already determined continuation value as another coordinate does not distinguish the two policy choices.
\end{example}

\subsection*{5) VERIFICATION}
The example uses the full distribution, not a moment approximation. Its failure is a policy-selection collision on a positive-mass individual state, not a null-set disagreement. All rewards and assets are bounded, so no transversality problem drives the failure. Linear utility is deliberate: a claim restricted to strictly concave MA-01 preferences is not disproved by this auxiliary economy. The finite check in the script verifies that both pushforwards of $(1/2,1/2)$ are $(1/2,1/2)$ and that $5/2+1/2=3$.

In the theorem, compactness is not misused to imply measurable selection for arbitrary data; the needed Bellman measurability and maximizing selector are explicit assumptions. Invariance includes all positive-probability aggregate branches. With unbounded rewards, the proof would need convergence of the discounted utility series and
$\beta^N\E|V(X_N,S_N)|\to0$ for the equilibrium and every deviation used in verification. An Euler equation alone does not supply those limits.

\subsection*{6) FINAL}
\unresolved

\textbf{Strongest checked partial result.} On the declared bounded Markov specification, recursive equilibria induce sequential equilibria for every initial state in the invariant domain, and a given full-continuation sequential equilibrium admits a reduced-state representation under the explicit fibre and measurability conditions. The auxiliary economy proves that optimal continuation values can be state-dependent while selected optimal policies are not.

\textbf{Exact first gap for the original programme.} After the missing MA-01 primitives are fixed, prove that a full sequential equilibrium correspondence admits a Borel continuation selection whose prices and successor aggregate distributions depend only on $(z,\mu)$, and whose household policies are consistent across repeated states. Neither household Bellman contraction nor existence of some sequential equilibrium proves this selection assertion.

\textbf{Next three moves.} (1) Prove fibre consistency of MA-01 continuation prices and laws. (2) Establish measurable selection on an invariant domain covering the required initial distributions. (3) Verify transversality for the actual utility and borrowing constraints.

\textbf{Likely minimal obstruction.} Repeated aggregate states may support distinct continuation equilibria, not merely household ties. The two-atom example isolates this selection issue without disproving an unspecified stricter economy.

\clearpage
\begin{thebibliography}{99}
\bibitem{Cao}
Dan Cao. \emph{Recursive Equilibrium in Krusell and Smith (1998)}. Author working paper, title page dated September 2016, hosted in the AEA 2018 meeting materials. The version checked is the 72-page working paper, especially the abstract and introduction; the citation does not equate its expanded-state result with an unrestricted reduced-state theorem. \url{https://www.aeaweb.org/conference/2018/preliminary/paper/Ty859AnT}.

\end{thebibliography}
\end{document}
